La EDA en menores de 5 años es una serie mucho más regular y las diferencias entre modelos son pequeñas (Tabla~\ref{tab:eda}). La GRU obtiene el menor MAE de prueba (\nEGRUMAE{} episodios por provincia y semana) y el menor RMSE (\nEGRURMSE), con una mejora de \nredEPers\,\% sobre la persistencia y de \nredEArima\,\% sobre ARIMA, pero queda prácticamente empatada con el \textit{random forest} (\nERFMAE), que además tiene el menor error en validación y en escala logarítmica. Como en el dengue, la ventaja de las redes aparece en los horizontes largos: a una semana todos los modelos rondan los 10 episodios, mientras que a 8 semanas el \textit{Transformer} (\nEHOchoTr) reduce en \nredEHochoTrPers\,\% el error de la persistencia (\nEHOchoPers). El F1 es más bajo que en el dengue (cerca de 0,5) porque las semanas de alerta de EDA son subidas cortas sobre una estacionalidad estable; en esta tarea ARIMA, que modela mejor esa estacionalidad, es el que mejor anticipa las alertas.
